% IMT221_lab_report_template.tex
% Plantilla de Informe de Laboratorio - Estilo IEEE
% Curso: IMT-221 — Circuitos Electrónicos II
% Ingeniería Mecatrónica - Universidad Católica Boliviana "San Pablo"

\documentclass[conference]{IEEEtran}

% =======================================================
% PAQUETES RECOMENDADOS
% =======================================================
\usepackage[utf8]{inputenc}
\usepackage[spanish, es-tabla, es-nodecimaldot]{babel} % Español, usa 'Tabla' en lugar de 'Cuadro'
\usepackage{amsmath}
\usepackage{amssymb}
\usepackage{graphicx}
\usepackage{booktabs}
\usepackage{siunitx}
\usepackage{cite}
\usepackage{url}
\usepackage{array}
\usepackage{xcolor}
\usepackage[colorlinks=true, linkcolor=black, citecolor=black, urlcolor=black]{hyperref}

% Configuración básica de siunitx para el español
\sisetup{output-decimal-marker = {.}}

\begin{document}

% =======================================================
% BLOQUE DE TÍTULO Y AUTORES
% =======================================================
% Reemplace los datos entre llaves con la información real.
\title{Título del Laboratorio: Subtítulo Descriptivo y Específico}

\author{
    \IEEEauthorblockN{Nombre1 Apellido1, Nombre2 Apellido2, Nombre3 Apellido3}
    \IEEEauthorblockA{
        Ingeniería Mecatrónica \\
        Universidad Católica Boliviana ``San Pablo''\\
        Tarija, Bolivia \\
        correo1@ucb.edu.bo, correo2@ucb.edu.bo, correo3@ucb.edu.bo
    }
}

\maketitle

% =======================================================
% ABSTRACT (RESUMEN)
% =======================================================
\begin{abstract}
% Describa aquí en un solo párrafo (100-180 palabras):
% 1. El propósito del laboratorio.
% 2. El circuito o dispositivo estudiado.
% 3. El método utilizado.
% 4. El resultado cuantitativo principal (ej. reducción de error de X a Y).
% 5. La conclusión técnica más importante.
% NO convierta el abstract en una introducción extensa ni use citas bibliográficas aquí.
[Escriba su resumen aquí...]
\end{abstract}

% =======================================================
% KEYWORDS (PALABRAS CLAVE)
% =======================================================
\begin{IEEEkeywords}
% Incluya entre 3 y 5 términos en inglés o español.
operational amplifier, differential amplifier, input offset voltage, slew rate, gain-bandwidth product.
\end{IEEEkeywords}

% =======================================================
% I. INTRODUCCIÓN
% =======================================================
\section{Introducción}
% Responda de forma concisa:
% - ¿Qué fenómeno se estudia y por qué es relevante en ingeniería?
% - ¿Qué pregunta experimental aborda el laboratorio?
% - Termine la sección declarando claramente el objetivo del trabajo.
% - NO incluya resultados completos en esta sección.

[Escriba su introducción aquí...]

% =======================================================
% II. FUNDAMENTO TEÓRICO
% =======================================================
\section{Fundamento Teórico}
% Incluya únicamente la teoría necesaria para interpretar el experimento.

\subsection{Modelo del Circuito}
% Describa la topología base. Defina las variables de entrada y salida (V1, V2, Vo).
% Ejemplo de ecuación numerada y referenciada:
% La relación de salida ideal está dada por \eqref{eq:salida_ideal}.
%
% \begin{equation}
%     V_{o,\mathrm{ideal}} = k(V_2 - V_1)
%     \label{eq:salida_ideal}
% \end{equation}

\subsection{Parámetros Reales del Amplificador Operacional}
% Describa las no idealidades relevantes para su laboratorio (Ej. Lab 3: V_OS, I_B, I_OS, mismatch. Lab 4: GBW, SR).
% Toda especificación de un componente real debe provenir de su datasheet y citarse correctamente.

\subsection{Relaciones Dinámicas (Si aplica)}
% Si realiza mediciones de frecuencia o de Slew Rate (Lab 4), puede incluir ecuaciones como:
% \begin{equation}
%     f_{\mathrm{BW}} \approx \frac{\mathrm{GBW}}{A_N}
%     \label{eq:bw}
% \end{equation}
% \begin{equation}
%     SR_{\mathrm{req}} = 2\pi f V_p
%     \label{eq:sr}
% \end{equation}

% =======================================================
% III. METODOLOGÍA
% =======================================================
\section{Metodología}
% Permita que otra persona entienda exactamente cómo se realizó el experimento.

\subsection{Componentes e Instrumentos}
% Identifique los componentes físicos utilizados. No invente especificaciones.
\begin{table}[htbp]
    \centering
    \caption{Componentes e Instrumentación Utilizada}
    \label{tab:componentes}
    \begin{tabular}{llcl}
        \toprule
        \textbf{Elemento} & \textbf{Modelo / Valor} & \textbf{Cant.} & \textbf{Observación} \\
        \midrule
        Op-Amp & LM741 / TL074 & - & [Fabricante] \\
        Resistor & \SI{10}{\kilo\ohm} & - & Tolerancia $1\%$ \\
        Resistor & \SI{100}{\kilo\ohm} & - & Tolerancia $1\%$ \\
        Fuente DC & - & 1 & $\pm\SI{15}{\volt}$ / $\pm\SI{10}{\volt}$ \\
        Osciloscopio & - & 1 & - \\
        \bottomrule
    \end{tabular}
\end{table}

\subsection{Circuito Experimental}
% Inserte el esquema eléctrico. Descomente las siguientes líneas cuando tenga su figura.
% La plantilla compilará incluso si no incluye figuras inicialmente.
%
% \begin{figure}[htbp]
%     \centering
%     % \includegraphics[width=0.8\columnwidth]{figures/circuito.pdf}
%     \caption{Esquema del circuito experimental implementado (Reemplace con su figura).}
%     \label{fig:circuito}
% \end{figure}

\subsection{Procedimiento Experimental}
% Describa la secuencia REAL que ejecutaron.
% Ejemplo: 1) Verificación de pinout y medición de componentes, 2) Ajuste de alimentación, 3) Aplicación de señales, 4) Medición, 5) Troubleshooting, 6) Compensación.
% ADVERTENCIA: NO copie textualmente la guía del laboratorio. Escriba en pasado y en pasiva (ej. "Se midieron los resistores...").

\subsection{Simulación (Opcional o Complementaria)}
% Si utilizó LTSpice, documéntelo aquí.
% RECUERDE: Simulación NO es medición física. Indique el macromodelo utilizado.

% =======================================================
% IV. RESULTADOS
% =======================================================
\section{Resultados}
% Presente evidencia (datos medidos), sin interpretarla extensamente todavía.
% Reemplace los guiones (-) con sus mediciones. No fabrique datos.

\subsection{Medición DC (Lab 3)}
\begin{table}[htbp]
    \centering
    \caption{Comparación: Teoría vs Medición (Precisión DC)}
    \label{tab:resultados_dc}
    \begin{tabular}{ccccc}
        \toprule
        \textbf{Cond.} & \textbf{$V_{1}$} & \textbf{$V_{2}$} & \textbf{$V_{o,\mathrm{ideal}}$} & \textbf{$V_{o,\mathrm{meas}}$} \\
        \midrule
        T1 & \SI{0}{\volt} & \SI{0}{\volt} & \SI{0}{\volt} & - \\
        T2 & \SI{0}{\volt} & \SI{50}{\milli\volt} & \SI{0.5}{\volt} & - \\
        \bottomrule
    \end{tabular}
\end{table}

\subsection{Evaluación de Compensación}
\begin{table}[htbp]
    \centering
    \caption{Efecto de la Compensación de Offset (T1)}
    \label{tab:compensacion}
    \begin{tabular}{lcc}
        \toprule
        \textbf{Variable} & \textbf{Antes (Before)} & \textbf{Después (After)} \\
        \midrule
        $V_{o}$ medido & - & - \\
        Error absoluto ($|e|$) & - & - \\
        \bottomrule
    \end{tabular}
\end{table}

% \subsection{Mediciones Dinámicas (Para Lab 4)}
% \begin{table}[htbp]
%     \centering
%     \caption{Respuesta en Frecuencia y Slew Rate}
%     \label{tab:dinamico}
%     \begin{tabular}{cccc}
%         \toprule
%         \textbf{Frecuencia} & \textbf{$V_{\mathrm{in}}$} & \textbf{$V_{\mathrm{out}}$} & \textbf{Distorsión obs.} \\
%         \midrule
%         - & - & - & - \\
%         - & - & - & - \\
%         \bottomrule
%     \end{tabular}
% \end{table}

% =======================================================
% V. DISCUSIÓN
% =======================================================
\section{Discusión}
% ESTA ES LA SECCIÓN MÁS IMPORTANTE DEL INFORME.
% No repita los datos de la sección IV; interprételos.
% Responda:
% 1. ¿Coincidió el experimento con la predicción ideal?
% 2. ¿Qué discrepancia (error) apareció y qué mecanismos podrían explicarla?
% 3. ¿Qué evidencia permite descartar o apoyar cada hipótesis? (Ej. "Tras aplicar offset-null, el error se redujo en X mV, lo que sugiere que...")
% 4. ¿Qué diferencia hubo entre simulación y realidad?
% 5. Identifique limitaciones de sus mediciones.
%
% REGLA: Observación (Hecho) != Interpretación (Hipótesis) != Conclusión.

[Escriba su discusión aquí...]

% =======================================================
% VI. CONCLUSIONES
% =======================================================
\section{Conclusiones}
% Responda directamente al objetivo técnico planteado en la Introducción.
% Debe incluir el resultado cuantitativo principal y su significado técnico.
% Identifique la limitación principal o el siguiente paso lógico.
%
% ADVERTENCIA: NO escriba simplemente "El laboratorio fue exitoso" o "Aprendimos mucho".
% Ejemplo aceptable: "La compensación redujo el error DC de salida de X mV a Y mV. Sin embargo, persistió un error atribuible a las tolerancias resistivas..."

[Escriba sus conclusiones aquí...]

% =======================================================
% REFERENCIAS
% =======================================================
\begin{thebibliography}{00}
% Use el formato IEEE para citar. Sustituya estos ejemplos con sus fuentes reales.
% Los Datasheets SON fuentes primarias obligatorias.

\bibitem{ti_lm741}
Texas Instruments, ``LM741 Operational Amplifier,'' Datasheet, Oct. 2015. [Online]. Available: \url{http://www.ti.com}

\bibitem{ti_tl082}
Texas Instruments, ``TL08x JFET-Input Operational Amplifiers,'' Datasheet, May 2015. [Online]. Available: \url{http://www.ti.com}

\bibitem{textbook}
S. Franco, \emph{Design with Operational Amplifiers and Analog Integrated Circuits}, 4th ed. New York, NY, USA: McGraw-Hill, 2014.

\end{thebibliography}

% =======================================================
% APÉNDICE (OPCIONAL)
% =======================================================
% Descomente si necesita añadir datos crudos, cálculos muy extensos, 
% esquemas secundarios de LTSpice o código fuente auxiliar.
% 
% \appendices
% \section{Datos Crudos de Laboratorio}
% [Tablas extensas o configuraciones específicas...]
% 
% \section{Derivación Adicional}
% [Cálculos largos que interrumpirían la lectura del informe principal...]

\end{document}
